\chapter{理论基础与研究假设}

本章主要内容是建构研究框架并提出相应的假设。首先，系统论述计划行为理论、自我决定理论、社会认知理论的主要构成要素，进一步探究这三种理论对于创业意愿形成机理所具备的契合之处以及互动情况；其次，利用理论分析与实证检验相融合的方式，从教育模式对个体自主动机产生的影响、自主动机的中介效应对其内部运作的阐释作用、外部激励对这种效应起到的调节功能这三个方面提出具体的假设命题，并加以归纳整理；在上述理论推演和假设提出之后，就将各个理论相互推演的过程以及假设所包含的全部内容都用图形的形式呈现出来，从而清楚地体现出各个变量之间的逻辑关系以及整体结构。

\section{相关理论基础}

\subsection{计划行为理论}

计划行为理论是由Ajzen在1991年提出的，是社会科学领域里解释和预判人们做出行动意向的理论模型\cite{Ajzen1991}。该理论认为行为意向是评价个体实际行动的直接指标，行为意向是由行为态度、主观规范和感知行为控制力这三个主要前因变量所共同决定的。

在创业学的研究中，计划行为理论成为解释个体创业意向形成过程的主要分析模式。Al-Jubari等\nolinebreak[3]（2019）整合自我决定理论与计划行为理论，验证了态度、主观规范和知觉行为控制在自主动机与创业意向之间的中介作用\cite{AlJubari2019}。Kazmi和Nábrádi\nolinebreak[3]（2017）的研究也采用了计划行为理论作为解释创业教育效果的框架\cite{Kazmi2017}。武梦超和李礼旭\nolinebreak[3]（2025）的研究同样基于计划行为理论和社会认知理论的整合框架\cite{wumengchao2025}。本研究以计划行为理论为依据，对创业意愿的产生机理及其内在逻辑进行深入探究。

该理论认为行为意向是促使人去做某种行为的一种内部动因，它的形成是由诸多前因变量共同起作用的。行为态度是指一个人对某项行为所持有的主观喜好和价值取向，它是由该行为所可能产生的结果认知评价以及自身价值观的塑造所决定的。大学生如果认为参加创新创业活动有助于自身发展、取得经济收益或者获得社会认可的话，那么就会有较高的积极性去参与到创业当中来。主观规范是体现个人对个体化的社会期望的影响程度的指标，它是来自重要他人给出的信息传递，并且它引发的心理认同的机制。关于感知行为控制，它反映的是个体对完成某件事情所需要资源的掌控感、对所拥有的技能的自信，间接体现的是人们对自身未来行动可能性的心理预期。

计划行为理论在创业学的研究中得到了很多的实证检验和学术认可。Krueger等\nolinebreak[3]（2000）研究认为，创业者对目标的认识偏好、外在约束感知以及行为执行的主观信念都会直接决定创业意愿\cite{Krueger2000}。本研究选择计划行为理论为主导的分析框架，是因为它可以清楚地解释创业意愿形成的心理过程，在创新创业教育实践当中，通过促使学生形成积极的认知结构、增强职业认同感以及改善技能应用信心等途径来提高学生的参与动机。

\subsection{自我决定理论}

自我决定理论由Deci和Ryan于1985年提出，其核心在于探讨个体动机的内涵及其对心理发展的内在作用机制。动机分为自主型和外控型两种类型，自主型动机源于个人的价值取向和内在的兴趣，这样的内在驱动的学习行为更高效、稳定，外控型动机由外部的奖惩等环境因素决定，一般具有短暂的效果。

在创业教育的研究中，自我决定理论也被广泛地用来分析不同个体差异下创业课程的学习效果。Al-Jubari等\nolinebreak[3]（2019）整合自我决定理论与计划行为理论，发现自主需求、胜任需求和归属需求的满足能够显著提升创业意向\cite{AlJubari2019}。陈锦和张海锋\nolinebreak[3]（2025）发现，自主动机对学生的创业意愿具有显著影响，在创新创业学习中可对创业意愿及行为表现发挥积极作用\cite{chenjin2025}。这一发现为本研究探讨自主动机的中介作用提供了重要的理论基础。

自我决定理论主要论述了人类心理需求的本质特征，有三个基本部分，分别是自主性需求、能力归属感需求、社会联结感需求。自主性需求指个体对于行为自主权以及内在价值的认识倾向，即个人的行为是出于自己的意愿而不是被迫去做的；能力认同需求指的是个体在社会互动过程中对于自身效能的感知以及潜能开发的过程体验；关联性需求指人际交往的社会功能和情感纽带的意义，可以增强归属感并巩固身份认同。当这三种需要得到满足的时候，个体会激发起内在的动机，有很高的学习参与度，具有较强的创新思维以及情绪调控能力；而缺乏满足这些需求的时候，则会出现依赖性的动机或者消极的应对方式。

自我决定理论构建了动机发展的连续谱模型，对它的由外及内层次性展开系统分析。首层属于外部调节型动机，个体的行为主要受制于外在的奖惩，为了成绩或者避免负面评价而参加学习活动。第二层是内摄型动机，这时个体虽然还受外部规则的影响，但是已经产生了内部规范的意识，倾向于用规避负向情绪或者履行责任义务的方式来驱动自己的行为。第三层就是认同型动机，个体能认识到某项行为所具有的意义和价值，将其看成是达成自己核心目标的重要途径。最终形成的就是自主型动机，个体是依据活动的本质去理解，并且得到愉快的感受，从而无需外界的驱使而自己就会去做。

\subsection{社会认知理论}

社会认知理论是由班杜拉在1986年提出的，其主要观点是关于人、行为和环境三者之间相互影响的理论\cite{Bandura1986}。该理论中的核心概念就是\nolinebreak[3]“自我效能”，即个人对于达成某个目标的能力的主观评价水平。Lent等人\nolinebreak[3]（1994）在此基础上将该理论体系扩展到职业发展领域，并做了进一步的研究\cite{Lent1994}。

在创业学研究方面，自我效能感是决定创业者行为意向的最重要中介变量，已被大家所熟知并做进一步的研究。据已有文献显示，创业教育经由加强个人自我效能感这一途径，明显改善其创业动机状况。Bohlayer和Gielnik\nolinebreak[3]（2023）的研究关注到行动导向创业培训中创业自我效能感的下降现象，指出需要关注培训过程中可能出现的\nolinebreak[3]“训练困境”\cite{Bohlayer2023}。Gielnik等\nolinebreak[3]（2015）通过对学生创业培训的评估研究发现，行动调节在创业过程中发挥着关键作用\cite{Gielnik2015}。李其容等\nolinebreak[3]（2023）发现，创业自我效能在创业进展与创业努力之间发挥关键中介作用\cite{liqirong2023}。本研究用社会认知理论为依托，对自主学习型教育模式下创业意愿的影响路径和心理传导机制做了系统的分析。

社会认知理论的主要含义就体现在三元交互决定机制的阐释框架里，它研究的是个人特征、行为和环境三者之间的复杂互动。在此理论体系里，个人特质是由认知能力、情感状况、生理特性这些方面组成；具体的行为就是个人所表现出的行动表现；环境是指与行为的发生有关的各种情境因素以及条件限制。构成维度不是独立的线性作用链条，而是层层反馈、层层调节的过程，最终影响主体的发展轨迹。学生创业认知水平的高低会决定他是否会对创业机会展开探究，而实际体验会不断地修正和完善已有的认知。系统化的创业教育不但可以激发学生积极的创业意识，并且使学生付诸行动，而且还可以通过正向激励来加强师生之间的互动，从而改善整个教学环境。

自我效能感属于社会认知理论里的关键概念，它体现的是个体对于完成某项任务的能力信心程度。班杜拉认为自我效能感主要是通过成功经验、替代性经验、言语说服、生理唤醒和情绪状态这四个途径形成的。社会认知理论中一个核心的概念就是观察学习，观察学习是指个体通过观察别人的行为以及它们产生的结果，从而达到间接习得的一种方式。该理论给案例教学法、典型事迹讲授和实地考察这些多元化的实践教学方法赋予了关键的理论支撑。

\section{研究假设提出}

\subsection{教育模式与自主动机的关系}

按照自我决定理论，环境因素经由满足个体的自主性需求、能力认同感以及社会归属感，可促使内在动机得以形成并发展起来。创新创业教育课程属于重要的环境变量之一，创新创业教育课程设计的好坏会直接决定学生自主学习积极性被调动的程度。各种教育模式因为执行途径的不同，在促使学习者自主性动机方面存在明显的差别特点。

以系统化知识传授为特点的理论型课程，在激发学生创业意愿上起着不可替代的作用。Gielnik等\nolinebreak[3]（2015）的研究表明，行动导向的创业培训能够有效提升学生的创业行动调节能力\cite{Gielnik2015}。该类课程以建立完整的创业知识体系为出发点，一方面可以消除学生创业认知上的障碍，另一方面也能使学生对创业本质进行更深层次的理解。因此，本文假设如下：

\textbf{H1a：理论型创新创业教育模式正向影响自主动机。}

实践性课程以体验式学习与技能培养相结合为特点。Gielnik等\nolinebreak[3]（2017）从长期视角探讨了创业培训对激情的持续影响，发现实践导向的培训能够有效提升学生的创业热情和持久投入\cite{Gielnik2017}。该类课程通过模拟或者真实的情境来开展教学，使学生在实践中获得知识、锻炼能力，进而逐渐产生内在的自主学习动机。因此可以根据该理论框架提出如下假设：

\textbf{H1b：实践型创新创业教育模式正向影响自主动机。}

融合型课程是以理论和实践相结合为特点的课程，它重视认知领域预习性质的学习，并且也重视行为技能的系统训练。Manimala和Thomas\nolinebreak[3]（2017）认为综合性创业教育比目标定向型教育模式有更高的效用\cite{Manimala2017}。其设计原理是依靠科学的理论指导实践活动，用实践活动加深对理论内涵的理解，两者相互促进、互相支撑，从而有效地调动学生自主学习的积极性。

\textbf{H1c：融合型创新创业教育模式正向影响自主动机。}

\subsection{自主动机的中介作用}

自主动机作为个体行为的内在驱动力，在教育干预向行为意向转化的过程中扮演着关键的中介角色。根据自我决定理论，自主动机源于个体的内在兴趣或价值认同，与更好的学习效果、更强的行为持续性和更高的心理健康水平相关。

在创业教育情境中，课程质量作为环境因素，通过满足学生的基本心理需求来激发自主动机，而自主动机的提升又促进了创业意愿的形成。Al-Jubari等\nolinebreak[3]（2019）运用自我决定理论和计划行为理论的整合框架，发现自主动机通过态度、主观规范和知觉行为控制的中介作用影响创业意向\cite{AlJubari2019}。陈锦和张海锋\nolinebreak[3]（2025）的研究也证实了自主动机在创业学习与创业意愿之间的中介作用\cite{chenjin2025}。据此提出：

\textbf{H2a：自主动机在理论型课程与创业意愿之间起中介作用。}

\textbf{H2b：自主动机在实践型课程与创业意愿之间起中介作用。}

\textbf{H2c：自主动机在融合型课程与创业意愿之间起中介作用。}

\subsection{外部动机的调节作用}

外部动机作为由外部压力或奖励驱动的行为动力，在自我决定理论中被认为可能对内在动机产生\nolinebreak[3]“挤出效应”。在创业教育情境中，学生参与课程的动机呈现异质性，部分学生出于对创业的真正兴趣，属于自主动机驱动，而另一部分学生主要是为了获得学分、奖励或证书，属于外部动机驱动。当外部动机在学生的动机结构中占主导地位时，可能会削弱自主动机对创业意愿的促进作用。

Bohlayer和Gielnik\nolinebreak[3]（2023）的研究发现，行动导向的创业培训可能导致创业自我效能感的暂时性下降，这种现象在外部动机较高的学生中表现得更为明显\cite{Bohlayer2023}。这说明外部动机可能负向调节培训效果。与此同时，外部动机的调节作用可能因教育模式的不同而呈现差异。

在理论型课程中，学生即使动机不强，也可以通过被动听讲获取知识，外部奖励的诱惑相对容易满足，挤出效应可能相对较弱但仍然存在。据此提出假设H3a，外部动机负向调节理论型课程中自主动机与创业意愿的关系。但基于理论推导，提出：

\textbf{H3a：外部动机负向调节理论型课程中自主动机与创业意愿的关系。}

在实践型课程中，实践环节需要学生主动投入，如果学生仅仅是为了完成任务而参与，缺乏真正的兴趣，那么实践体验带来的自主动机提升可能难以转化为真实的创业意愿。据此提出假设H3b，外部动机负向调节实践型课程中自主动机与创业意愿的关系。据此提出：

\textbf{H3b：外部动机负向调节实践型课程中自主动机与创业意愿的关系。}

在融合型课程中，理论与实践的双重要求意味着学生需要在两个层面都保持较高投入，外部动机的工具性取向可能在这种情境下表现得最为明显，当学生将课程视为获取学分或证书的手段时，自主动机向创业意愿的转化受阻。据此提出假设H3c，外部动机负向调节融合型课程中自主动机与创业意愿的关系。据此提出：

\textbf{H3c：外部动机负向调节融合型课程中自主动机与创业意愿的关系。}

\section{理论模型}

根据前面的理论基础和研究假设，本研究得到如图\ref{fig:theoretical_model}所示的概念模型。该模型有三个主要要素，把创新创业教育模式\nolinebreak[3]（理论型、实践型、融合型三种类型）看作自变量，自主动机是重要的中介变量，创业意愿是因变量，外部动机是调节变量。基本逻辑是不同的创新创业教育模式会通过影响人的自主动机来影响人的创业意愿，存在明显的中介效应。此时外部动机起着重要的调节作用。

\begin{figure}[H]
    \centering
    \begin{tikzpicture}[
        node distance=0.35cm and 1.6cm,
        box/.style={draw, rectangle, minimum width=1.8cm, minimum height=0.7cm, font=\zihao{5}, inner sep=3pt},
        arr/.style={-{Stealth[length=2.2mm]}, thick}
    ]
        % 左侧三类课程
        \node[box] (theory) {理论型};
        \node[box, below=of theory] (practice) {实践型};
        \node[box, below=of practice] (fusion) {融合型};
        % 外框（虚线）
        \node[draw, dashed, fit=(theory)(practice)(fusion), inner sep=10pt, inner ysep=14pt] (modes) {};
        \node[font=\zihao{5}, anchor=south west, inner sep=2pt] at (modes.north west) {创新创业教育模式};
        % 中介与因变量
        \node[box, right=of modes, minimum width=1.6cm] (auto) {自主动机};
        \node[box, right=of auto, minimum width=1.6cm] (intent) {创业意愿};
        % 调节变量
        \node[box, above=1.2cm of $(auto.east)!0.5!(intent.west)$, minimum width=1.6cm] (ext) {外部动机};
        % 箭头
        \draw[arr] (modes.east) -- (auto.west);
        \draw[arr] (auto.east) -- (intent.west);
        \draw[arr] (ext.south) -- ($(auto.east)!0.5!(intent.west)$);
    \end{tikzpicture}
    \caption{理论模型图}
    \label{fig:theoretical_model}
\end{figure}

本研究以计划行为理论、自我决定理论和社会认知理论为基础，对创业意愿的形成机理做系统的探究。根据计划行为理论框架，创业态度和感知控制力一起成为影响个体创业意图的主要因素，从自我决定理论视角看，自主性动机既是外部激励效能的中介变量，又是行为动机的调节变量；根据社会认知理论分析路径，在教育干预促进创业意愿提高过程中，自主性动机起到关键的中介作用。为清晰呈现本研究提出的各项假设，现将所有假设汇总于表\ref{tab:hypothesis_summary}。

\tablemacro{!htb}
            {假设汇总表}
            {cp{9.5cm}c}
            {\songtibf{假设编号} & \songtibf{假设内容} & \songtibf{预期方向} \\}
            {H1a & 理论型创新创业教育模式正向影响自主动机 & 正向 \\
            H1b & 实践型创新创业教育模式正向影响自主动机 & 正向 \\
            H1c & 融合型创新创业教育模式正向影响自主动机 & 正向 \\
            H2a & 自主动机中介理论型课程与创业意愿 & --- \\
            H2b & 自主动机中介实践型课程与创业意愿 & --- \\
            H2c & 自主动机中介融合型课程与创业意愿 & --- \\
            H3a & 外部动机负向调节理论型课程中自主动机与创业意愿 & 负向 \\
            H3b & 外部动机负向调节实践型课程中自主动机与创业意愿 & 负向 \\
            H3c & 外部动机负向调节融合型课程中自主动机与创业意愿 & 负向 \\}
            {tab:hypothesis_summary}

\clearpage
